% !PRINTED 238
\weberpage{238}
%<<CONT
one recognises that $j(\omega)$ can be developed into a series
progressing according to powers of $q^2$, which has the form
\weq{4}{j(\omega) = q^{-2} + a_1 + a_2 q^2 + a_3 q^4 + \cdots,}
in which the $a_1$, $a_2$, $a_3 \ldots$ are \gesperrt{integers} which
can be calculated successively (there results, for example, $a_1 = 744$,
$a_2 = 196\,884$). The developments of the quantities (1) therefore
begin with
\weq{5}{e^{-\frac{2\pi i c}{a}} e^{-\frac{2\pi i \partial}{a}\omega}}
\gesperrt{and are therefore all different from one another. The
invariant equation is therefore, by § 65, irreducible.}

There is, however, a second way of arriving at this equation, and at
the transformation equations generally, which we shall now first
become acquainted with in the case of the invariant equation.
This derivation rests upon the theorems of § 54 on modular functions
and holds also for an \gesperrt{even} $n$. Here it is
Theorem 4, § 54, that comes into application:

I. \gesperrt{Every modular function which under the two transformations}
\weq{6}{\begin{pmatrix} 1, & 0 \\ 1, & 1 \end{pmatrix},}
\weq{7}{\begin{pmatrix} 0, & 1 \\ -1, & 0 \end{pmatrix}}
\gesperrt{remains unchanged is a rational function of $j(\omega)$.}

We first show that by application of the substitutions
(6), (7) the $\nu$ quantities (1) are permuted among one another.
We have the composition
\weq{8}{\begin{pmatrix} a, & 0 \\ c, & \partial \end{pmatrix}\begin{pmatrix} 1, & 0 \\ 1, & 1 \end{pmatrix} = \begin{pmatrix} 1, & 0 \\ \lambda, & 1 \end{pmatrix}\begin{pmatrix} a, & 0 \\ c_1, & \partial \end{pmatrix},}
when
\weq{9}{c_1 = c + \partial - \lambda a}
is set. Since $j(\omega)$ remains unchanged under every linear
transformation, transformation (6) acts as follows:
\[
j\left(\frac{c+\partial\omega}{a}\right) \text{ goes to } j\left(\frac{c_1+\partial\omega}{a}\right)
\]

Further we determine $a_2$, $\partial_2$, $c_2$ so that
\weq{10}{\begin{gathered}
\begin{pmatrix} a, & 0 \\ c, & \partial \end{pmatrix}\begin{pmatrix} 0, & 1 \\ -1, & 0 \end{pmatrix} = \begin{pmatrix} \alpha, & \beta \\ \gamma, & \delta \end{pmatrix}\begin{pmatrix} a_2, & 0 \\ c_2, & \partial_2 \end{pmatrix}, \\
\alpha\delta - \beta\gamma = 1.
\end{gathered}}
